\section{Focused Real/Pseudo-Real Benchmark}
Table~\ref{tab:focused-real} shows the main focused benchmark under structured contamination. The clean target rank is the last rank in each dataset; preserving that rank means the attack was blocked.

\begin{table}[H]
\centering
\small
\caption{Focused real/pseudo-real target-rank results under structured contamination.}
\label{tab:focused-real}
\begin{tabular}{lrrrrrrrrr}
\toprule
Dataset & Clean & M1 & M2 & M3 & M4 & M5 & M6 & M7 \\
\midrule
Healthcare countries 2021 & 26 & 1 & 7 & 26 & 26 & 5 & 26 & 26 \\
Car evaluation & 300 & 24 & 180 & 300 & 300 & 192 & 300 & 300 \\
Healthcare resource allocation & 300 & 1 & 52 & 300 & 300 & 16 & 300 & 300 \\
\bottomrule
\end{tabular}
\end{table}


M1 collapses under contamination. M2 improves over M1 but remains insufficient. M4, M6, and M7 preserve the clean target rank in all three focused larger datasets. M5 is inconsistent.

\section{Attack-Fraction Curves}
Table~\ref{tab:attack-curve} shows target-rank behavior as the effective attack fraction increases. M4 and M6 remain strong through 40\% effective structured contamination. M7 preserves the clean target rank across all tested structured contamination levels, including 53.3\% and 60.0\%.

\begin{longtable}{llrrrrrrrr}
\caption{Target-rank behavior across attack fractions. Clean rank is the desired preserved target rank.}
\label{tab:attack-curve}\\
\toprule
Dataset & Eff. & Clean & M1 & M2 & M3 & M4 & M5 & M6 & M7 \\
\midrule
\endfirsthead
\toprule
Dataset & Eff. & Clean & M1 & M2 & M3 & M4 & M5 & M6 & M7 \\
\midrule
\endhead
Healthcare countries & 13.3\% & 26 & 1 & 18 & 26 & 26 & 19 & 26 & 26 \\
Healthcare countries & 20.0\% & 26 & 1 & 11 & 26 & 26 & 11 & 26 & 26 \\
Healthcare countries & 26.7\% & 26 & 1 & 7 & 26 & 26 & 5 & 26 & 26 \\
Healthcare countries & 40.0\% & 26 & 1 & 2 & 26 & 26 & 26 & 26 & 26 \\
Healthcare countries & 53.3\% & 26 & 1 & 1 & 1 & 1 & 1 & 1 & 26 \\
Healthcare countries & 60.0\% & 26 & 1 & 1 & 1 & 1 & 1 & 1 & 26 \\
Car evaluation & 13.3\% & 300 & 78 & 272 & 300 & 300 & 289 & 300 & 300 \\
Car evaluation & 20.0\% & 300 & 48 & 230 & 300 & 300 & 258 & 300 & 300 \\
Car evaluation & 26.7\% & 300 & 24 & 180 & 300 & 300 & 192 & 300 & 300 \\
Car evaluation & 40.0\% & 300 & 1 & 77 & 300 & 300 & 300 & 300 & 300 \\
Car evaluation & 53.3\% & 300 & 1 & 18 & 1 & 1 & 1 & 1 & 300 \\
Car evaluation & 60.0\% & 300 & 1 & 8 & 1 & 1 & 1 & 1 & 300 \\
Healthcare allocation & 13.3\% & 300 & 3 & 168 & 300 & 300 & 166 & 300 & 300 \\
Healthcare allocation & 20.0\% & 300 & 1 & 101 & 300 & 300 & 62 & 300 & 300 \\
Healthcare allocation & 26.7\% & 300 & 1 & 52 & 300 & 300 & 16 & 300 & 300 \\
Healthcare allocation & 40.0\% & 300 & 1 & 11 & 300 & 300 & 300 & 300 & 300 \\
Healthcare allocation & 53.3\% & 300 & 1 & 1 & 1 & 1 & 1 & 1 & 300 \\
Healthcare allocation & 60.0\% & 300 & 1 & 1 & 1 & 1 & 1 & 1 & 300 \\
\bottomrule
\end{longtable}


\section{External Comparator Baselines}
Table~\ref{tab:external-comparators} compares the proposed methods with external comparator baselines. Robust aggregation and consensus filtering work well at low-to-moderate contamination but collapse under high contamination. The Huang--Li group-ideal adaptation is vulnerable because the original model assumes honest decision makers and does not include adversarial reliability detection.

\begin{longtable}{llrrrrrr}
\caption{External comparator target-rank results.}
\label{tab:external-comparators}\\
\toprule
Dataset & Attack & Clean & Median & Trimmed & MAD & Borda & Huang--Li \\
\midrule
\endfirsthead
\toprule
Dataset & Attack & Clean & Median & Trimmed & MAD & Borda & Huang--Li \\
\midrule
\endhead
Healthcare countries & 10\% & 26 & 26 & 26 & 26 & 25 & 8 \\
Healthcare countries & 20\% & 26 & 26 & 26 & 26 & 22 & 6 \\
Healthcare countries & 30\% & 26 & 26 & 19 & 26 & 22 & 4 \\
Healthcare countries & 40\% & 26 & 23 & 1 & 26 & 16 & 2 \\
Healthcare countries & 50\% & 26 & 1 & 1 & 1 & 12 & 1 \\
Healthcare countries & 60\% & 26 & 1 & 1 & 1 & 9 & 1 \\
Car evaluation & 10\% & 300 & 300 & 300 & 300 & 290 & 118 \\
Car evaluation & 20\% & 300 & 300 & 300 & 300 & 282 & 73 \\
Car evaluation & 30\% & 300 & 300 & 293 & 300 & 267 & 53 \\
Car evaluation & 40\% & 300 & 300 & 98 & 300 & 229 & 27 \\
Car evaluation & 50\% & 300 & 1 & 1 & 1 & 121 & 15 \\
Car evaluation & 60\% & 300 & 1 & 1 & 1 & 85 & 11 \\
Healthcare allocation & 10\% & 300 & 300 & 300 & 300 & 275 & 17 \\
Healthcare allocation & 20\% & 300 & 299 & 300 & 300 & 257 & 12 \\
Healthcare allocation & 30\% & 300 & 292 & 170 & 300 & 238 & 8 \\
Healthcare allocation & 40\% & 300 & 278 & 1 & 300 & 179 & 5 \\
Healthcare allocation & 50\% & 300 & 1 & 1 & 1 & 127 & 4 \\
Healthcare allocation & 60\% & 300 & 1 & 1 & 1 & 95 & 3 \\
\bottomrule
\end{longtable}


\section{Statistical Summary}
Table~\ref{tab:method-summary} summarizes blocked cases and target-rank error across 18 attack-fraction scenarios. M7 is the only method with 18/18 blocked cases and zero mean target-rank error under the tested structured attack model.

\begin{table}[H]
\centering
\small
\caption{Target-rank-error summary across 18 attack-fraction scenarios.}
\label{tab:method-summary}
\begin{tabular}{lrrrr}
\toprule
Method & Cases & Blocked & Rate & Mean error \\
\midrule
EB1 Median TOPSIS & 18 & 8 & 0.444 & 71.11 \\
EB2 Trimmed-Mean TOPSIS & 18 & 6 & 0.333 & 106.44 \\
EB3 MAD-Consensus TOPSIS & 18 & 12 & 0.667 & 69.22 \\
EB4 Individual-Borda TOPSIS & 18 & 0 & 0.000 & 66.94 \\
EB5 Huang--Li Group-Ideal TOPSIS & 18 & 0 & 0.000 & 188.22 \\
M1 & 18 & 0 & 0.000 & 199.39 \\
M2 & 18 & 0 & 0.000 & 144.95 \\
M3 & 18 & 12 & 0.667 & 69.22 \\
M4 & 18 & 12 & 0.667 & 69.22 \\
M5 & 18 & 3 & 0.167 & 117.03 \\
M6 & 18 & 12 & 0.667 & 69.22 \\
M7 & 18 & 18 & 1.000 & 0.00 \\
\bottomrule
\end{tabular}
\end{table}


\section{Pairwise Target-Error Tests}
M7 had lower target-rank error than M1, M2, EB4, and EB5 in all 18 attack-fraction scenarios. Compared with M3, M4, M6, and EB3, M7 improved the six majority/high-contamination scenarios and tied the remaining twelve. The sign-test and Wilcoxon-style paired comparisons support the improvement of M7 under the structured threat model.

\begin{table}[H]
\centering
\small
\caption{Selected pairwise target-error tests comparing M7 against major baselines.}
\label{tab:pairwise-tests}
\begin{tabular}{lrrrr}
\toprule
Comparison & M7 better & Ties & Sign-test \(p\) & Wilcoxon \(p\) \\
\midrule
M7 vs M1 & 18 & 0 & 0.000008 & 0.000214 \\
M7 vs M2 & 18 & 0 & 0.000008 & 0.000214 \\
M7 vs M4 & 6 & 12 & 0.031250 & 0.036032 \\
M7 vs M6 & 6 & 12 & 0.031250 & 0.036032 \\
M7 vs EB5 Huang--Li & 18 & 0 & 0.000008 & 0.000214 \\
\bottomrule
\end{tabular}
\end{table}


\section{Runtime and Scalability}
Table~\ref{tab:runtime} reports representative runtime at the 30\% contaminated setting with 200 bags. M7 is slower than M1 but comparable to the other ensemble methods and practical for offline decision-support use.

\begin{table}[H]
\centering
\small
\caption{Representative one-run runtime at 30\% contamination with 200 bags.}
\label{tab:runtime}
\begin{tabular}{lrrrrr}
\toprule
Dataset & M1 & M2 & M4 & M6 & M7 \\
\midrule
Healthcare countries & 0.004 & 0.380 & 0.387 & 0.291 & 0.315 \\
Car evaluation & 0.028 & 2.196 & 2.226 & 1.728 & 1.818 \\
Healthcare allocation & 0.077 & 6.809 & 6.863 & 5.468 & 5.731 \\
\bottomrule
\end{tabular}
\end{table}


\section{Hand-Computed Validation}
A small hand-computed fuzzy TOPSIS example was used as a smoke test. The expected closeness coefficients are \(CC(A_1)=0.52941176\) and \(CC(A_2)=0.47058824\), which match the implementation output.

\section{Result Interpretation by Research Question}

The results answer the research questions as follows.

\textbf{RQ1: Is classical fuzzy TOPSIS vulnerable?} Yes. M1 repeatedly moves the attacked target upward, often to rank 1. This confirms that standard fuzzy TOPSIS is not reliable under coordinated contamination.

\textbf{RQ2: Does bootstrap bagging alone solve the problem?} No. M2 improves over M1 in some cases but never fully blocks the structured target attack across the attack-fraction scenarios. Bagging reduces variance, but it does not identify biased decision makers.

\textbf{RQ3: Does reliability-weighted sampling help?} Yes. M4 preserves the clean target rank through 40\% effective contamination in the main attack-fraction experiments.

\textbf{RQ4: Does reliability-weighted aggregation help?} Yes. M6 also preserves the clean target rank through 40\% effective contamination and is easier to interpret than M7.

\textbf{RQ5: Does multi-signal reliability help under high contamination?} Yes. M7 is the only method that preserves the clean target rank in all 18 structured attack-fraction scenarios.

\textbf{RQ6: Are external comparator baselines enough?} The external baselines are useful but not enough under high contamination. Median, trimmed mean, and MAD-consensus methods can resist low or moderate contamination, but they collapse when the contaminated group shifts the center. The Huang--Li group-ideal adaptation is important prior art, but it assumes honest decision makers and is vulnerable under adversarial contamination.

\section{Frozen Result Decision}

The result tables in this chapter are frozen for the thesis report. The selected interpretation is:

\begin{quote}
M4, M6, and M7 are the three proposed methods. M2 is the corrected bootstrap baseline, and M5 is a cluster-stratified ablation branch. M7 is the strongest final method under the tested structured threat model.
\end{quote}
